% ============================================================
%  Chapter 8: Conclusions
%  Henry Kabuye PhD Thesis — Northumbria University
% ============================================================

\chapter{Conclusions}
\label{chap:conclusions}

\lettrine{T}{his} thesis has investigated the design, implementation, and evaluation of
a multi-stage, uncertainty-aware, and user-centric framework for automated ransomware
attack detection.
Beginning from the limitations of existing single-modality, overconfident, and opaque
detection systems, the research has progressed through three empirical studies — each
building on the verified contributions of its predecessor — to produce the Extended
Agentic SABRE framework: a three-agent neuro-symbolic architecture that integrates
statistical detection, behavioural analysis, and threat-intelligence reasoning with
calibrated uncertainty quantification, analyst-interpretable explanations, and
continual learning from expert feedback.

% ============================================================
\section{Summary of Contributions}
\label{sec:conc_contributions}
% ============================================================

% TODO: List the principal contributions, one per paragraph, tied to the
% specific results that evidence them. Recommended structure:
%
% Contribution 1 (Chapter 3): Single-agent uncertainty-aware detection baseline.
%  - MC Dropout uncertainty for ransomware CNN
%  - CTGAN-augmented class-imbalance handling
%  - SHAP/LIME analyst-interpretable explanations
%  - Results: F1/AUC figures on RanSAP and RDset
%
% Contribution 2 (Chapter 4): Two-agent Agentic SABRE with triage policy.
%  - Modality-specialised Semantic and Behavioural Agents
%  - Learned Decision MLP with score-space CTGAN augmentation
%  - Temperature-calibrated agent posteriors
%  - Bayesian TPE triage threshold optimisation
%  - Counterfactual explanation protocol
%  - Results: Semantic acc/AUC, Behavioural acc/AUC, triage escalation rate
%
% Contribution 3 (Chapter 5): Extended Agentic SABRE — three agents, peer
%  consultation, continual learning.
%  - RansKnow threat-intelligence dataset (307 samples, 23 features)
%  - Knowledge Agent with adaptive MC Dropout
%  - Peer Consultation Protocol (inverse-uncertainty-weighted fusion)
%  - Expert Feedback Loop with reservoir-sampled experience replay
%  - Results: 3-agent fusion performance, 9% escalation rate, continual
%    learning accuracy retention figures

% ============================================================
\section{Research Questions Answered}
\label{sec:conc_rqs}
% ============================================================

% TODO: Return to each numbered research question stated in Chapter 1 and
% provide a direct, evidenced answer drawn from the empirical results.
% Format: RQ → Answer (2--3 sentences) → Key supporting result.

% ============================================================
\section{Future Research Directions}
\label{sec:conc_future}
% ============================================================

% TODO: Identify the most important open problems that follow from this work.
% Suggested directions:
%
% 1. Live deployment evaluation: Deploy the framework on an institutional
%    endpoint monitoring system and evaluate real-world escalation rates,
%    analyst acceptance, and detection latency under production conditions.
%
% 2. Adversarial robustness: Evaluate the framework against evasion attacks
%    (feature-space adversarial examples) and investigate adversarially robust
%    training for the CNN backbone agents.
%
% 3. Richer knowledge representation: Replace the tabular RansKnow features
%    with a knowledge-graph representation of ATT&CK technique relationships,
%    enabling the Knowledge Agent to reason over tactic co-occurrence structure.
%
% 4. LLM-augmented threat intelligence: Explore the use of large language model
%    embeddings of practitioner threat reports as additional Knowledge Agent
%    input features.
%
% 5. Formal analyst user study: Conduct a user study with security analysts
%    to evaluate the utility of counterfactual and SHAP explanations in
%    supporting triage decisions, and to measure the quality of labels fed
%    back through the Expert Feedback Loop.
%
% 6. Extension to adjacent threats: Apply the multi-agent triage framework
%    to adjacent threat categories (banking trojans, APT lateral movement)
%    where the same multi-modal signal decomposition is applicable.

% ============================================================
\section{Closing Statement}
\label{sec:conc_closing}
% ============================================================

% TODO: A short closing paragraph (5--8 sentences) that restates the core
% thesis in plain language — what was the problem, what was built, and what
% does it demonstrate about the feasibility and design principles of
% uncertainty-aware multi-agent detection.
% End on a forward-looking note about what the framework enables.
